\chapter{Giới thiệu}
\label{chap:intro}

\section{Bối cảnh và động lực}

Việt Nam đang già hóa dân số với tốc độ nhanh. Theo số liệu năm 2024 của Tổng cục Thống kê, cả
nước có khoảng 14,2 triệu người cao tuổi, chiếm khoảng 14\% dân số, và dự kiến bước vào giai
đoạn dân số già vào khoảng năm 2034--2036~\cite{vnexpress2024tuoigia}. Việc chăm sóc người cao
tuổi hiện chủ yếu do người thân trong gia đình đảm nhận, trong khi dịch vụ chăm sóc cộng đồng
còn hạn chế. Bối cảnh đó tạo ra nhu cầu về những công cụ giúp người thân theo dõi sinh hoạt
hằng ngày của người cao tuổi từ xa, với chi phí thấp và không xâm phạm quyền riêng tư.

Nhà thông minh dùng \textbf{cảm biến môi trường} --- cảm biến chuyển động hồng ngoại (PIR), cảm
biến cửa, cảm biến nhiệt độ --- là một hướng đáp ứng nhu cầu đó: thiết bị rẻ, không cần camera,
không bắt người dùng đeo thiết bị~\cite{cook2013casas,bouchabou2021survey}. Từ nhật ký sự kiện
của các cảm biến, hệ thống có thể học được \textbf{thói quen} của người ở: thức dậy lúc mấy giờ,
ăn sáng khi nào, có uống thuốc đều đặn không. Thói quen là nền tảng cho nhiều ứng dụng: nhắc nhở
khi quên uống thuốc --- một vấn đề có thật khi tỉ lệ tuân thủ điều trị bệnh mạn tính trung bình
chỉ khoảng 50\%~\cite{who2003adherence} ---, cảnh báo khi sinh hoạt thay đổi bất thường, hay tự
động hóa thiết bị theo nhịp sống của từng người.

Trong đợt thực tập trước đó, tác giả đã xây dựng một phương pháp không giám sát để tìm thói quen
từ dữ liệu cảm biến: gom sự kiện thành \emph{phiên} (một lần ở liên tục trong một phòng), biểu
diễn giờ trong ngày trên vòng tròn và phân cụm bằng DBSCAN~\cite{ester1996dbscan}. Phương pháp
tìm lại đủ 12/12 thói quen cài sẵn trên dữ liệu mô phỏng và 5 thói quen trên dữ liệu thật của
một người sống một mình. Tuy nhiên, phương pháp đó --- cũng như phần lớn nghiên cứu về khai phá
mẫu sinh hoạt --- ngầm giả định \textbf{trong nhà chỉ có một người}.

Giả định này không đúng với thực tế. Người cao tuổi thường sống cùng vợ hoặc chồng, hoặc cùng con
cháu. Khi trong nhà có hai người, cảm biến môi trường \emph{không biết ai đã kích hoạt nó}: các
sự kiện của hai người trộn vào một luồng duy nhất~\cite{benmansour2015survey,wang2021smrt}. Luận
án tiến sĩ của Dong Chen~\cite{chen2021thesis} --- tài liệu định hướng chính của khóa luận ---
đã chỉ ra rằng nhận dạng hoạt động trong nhà nhiều người khó hơn hẳn nhà một người, và đề xuất
các kỹ thuật phân đoạn, phân lớp và học trực tuyến cho bài toán này. Nhưng luận án đó dừng ở việc
nhận dạng \emph{hoạt động đang diễn ra}; chính tác giả nêu trong phần hạn chế rằng các phương
pháp ``chỉ nhận ra hoạt động hiện tại'' và chưa hướng tới hành vi dài hạn.

\section{Bài toán và ví dụ minh họa}
\label{sec:intro-example}

Xét bộ dữ liệu CASAS Cairo~\cite{casas2025zenodo}: một cặp vợ chồng (ký hiệu R1, R2) sống trong
căn nhà ba tầng gắn 27 cảm biến chuyển động, được ghi nhận trong 57 ngày. Bảng~\ref{tab:cairo-morning}
thống kê bốn hoạt động buổi sáng.

\begin{table}[htbp]
\centering
\caption{Bốn hoạt động buổi sáng ở nhà Cairo (thống kê trên nhãn thật của toàn bộ dữ liệu)}
\label{tab:cairo-morning}
\begin{tabular}{llcl}
\toprule
Hoạt động & Người & Giờ trung vị & Cảm biến kích hoạt nhiều nhất \\
\midrule
Thức dậy & R1 & 06:44 & M005, M001 (tầng trên) \\
Thức dậy & R2 & 07:04 & M005, M009 (tầng trên) \\
\textbf{Uống thuốc} & \textbf{R2} & \textbf{07:20} & \textbf{M021, M024 (bếp)} \\
Ăn sáng & chung & 07:29 & M024, M021 (bếp) \\
\bottomrule
\end{tabular}
\end{table}

Thói quen uống thuốc của R2 và bữa sáng chung dùng đúng những cảm biến bếp giống nhau và chỉ cách
nhau khoảng 10 phút. Một phương pháp chỉ nhìn ``phòng nào, mấy giờ'' như ở đợt thực tập không thể
tách hai thói quen này; thực nghiệm ở Chương~\ref{chap:exp} cho thấy nó còn gộp mọi bữa ăn trong
ngày của vùng bếp vào \emph{một} cụm duy nhất. Muốn biết ``R2 có uống thuốc lúc khoảng 7 giờ mỗi
sáng không'', hệ thống phải (i) gán được các đoạn sự kiện cho đúng người và đúng hoạt động, rồi
(ii) từ đó rút ra thói quen của từng người.

Khóa luận giải bài toán: \textbf{từ luồng sự kiện cảm biến trộn lẫn của nhà hai người, tìm đúng
thói quen của từng người}, và chứng minh thói quen tìm được dùng được để \textbf{phát hiện ngày bỏ
thuốc}.

\section{Câu hỏi nghiên cứu}

Khóa luận trả lời ba câu hỏi nghiên cứu:

\begin{description}[leftmargin=1.3cm,style=sameline]
  \item[RQ1] Gán các đoạn sự kiện cảm biến cho \emph{đúng người và đúng hoạt động} chính xác đến
  đâu, khi chỉ dùng cảm biến môi trường và đánh giá theo thứ tự thời gian?
  \item[RQ2] Từ kết quả của RQ1, thói quen của từng người tìm được \emph{đúng đến đâu} so với thói
  quen đáp án? Lỗi ở khâu nhận dạng --- bỏ sót, lệch giờ, nhận nhầm, \emph{nhầm người} --- lan
  truyền sang thói quen như thế nào? (Câu hỏi trung tâm.)
  \item[RQ3] Thói quen học được có phát hiện được \emph{ngày bỏ thuốc} không, và tốt hơn cách tiếp
  cận theo phòng đến đâu?
\end{description}

\section{Mục tiêu và phạm vi}

\textbf{Mục tiêu.} Xây dựng và đánh giá một quy trình đầu--cuối đi từ nhật ký cảm biến thô tới
danh sách thói quen của từng người; đưa ra định nghĩa đo được của khái niệm ``thói quen theo từng
người'' cùng bộ chỉ số đánh giá ở mức thói quen.

\textbf{Phạm vi} (được chốt trước khi thực hiện để tránh thay đổi giữa chừng):
\begin{itemize}
  \item \emph{Trong phạm vi:} nhà hai người; dữ liệu công khai CASAS (Cairo, Tulum1, Tulum2 là nhà
  hai người; Milan, Aruba là nhà một người dùng để đối chứng); đánh giá ngoại tuyến, chia dữ liệu
  theo ngày; các phương pháp baseline và mô hình học trình bày ở Chương~\ref{chap:method}.
  \item \emph{Ngoài phạm vi:} camera, thiết bị đeo; mô hình ngôn ngữ lớn và mô hình nền; nhà từ ba
  người trở lên; tự thu thập dữ liệu; gán người cho \emph{từng} sự kiện cảm biến (dữ liệu không
  có nhãn này); học trực tuyến; dự đoán hành động tiếp theo; ứng dụng nhắc nhở thật, điều khiển
  thiết bị, robot.
\end{itemize}

\section{Đóng góp}

\begin{enumerate}
  \item \textbf{Định nghĩa và thuật toán thói quen.} Phát biểu hình thức khái niệm thói quen theo
  từng người và hiện thực nó thành thuật toán \emph{quét cửa sổ} đếm theo ngày, khắc phục hiện
  tượng ``chuỗi hóa'' của DBSCAN mà đề tài thực tập gặp phải.
  \item \textbf{Bộ đánh giá ở mức thói quen.} Thói quen đáp án, phép so khớp một--một bằng thuật
  toán Hungary và các chỉ số precision, recall, F1, sai số giờ, sai số tần suất; tách riêng đánh
  giá thói quen gắn với từng người.
  \item \textbf{So sánh có kiểm soát} năm cách tiếp cận (hai baseline không nhãn, một baseline tra
  bảng, Random Forest, Bi-GRU) và giới hạn trên Oracle trên năm nhà, đánh giá theo thứ tự thời gian.
  \item \textbf{Phân tích lỗi lan truyền}: làm nhiễu nhãn thật có kiểm soát để đo thói quen nhạy
  với từng loại lỗi nhận dạng đến đâu, kể cả lỗi \emph{tráo người} đặc trưng của nhà nhiều người.
  \item \textbf{Thí nghiệm ứng dụng} phát hiện ngày bỏ thuốc trên dữ liệu thật có cài bất thường.
  \item Mã nguồn mở, có kiểm thử, tái lập toàn bộ số liệu bằng một lệnh.
\end{enumerate}

\section{Ý nghĩa khoa học và thực tiễn}

\textbf{Về khoa học}, khóa luận nối hai hướng nghiên cứu đang tách rời: nhận dạng hoạt động trong
nhà nhiều người~\cite{chen2021thesis,wang2021smrt,chen2026mdp} và khai phá mẫu sinh hoạt vốn mới
được kiểm chứng trên nhà một người~\cite{karpekov2025discover}. Việc đo lỗi lan truyền cho biết
khâu nhận dạng cần tốt đến đâu thì thói quen mới đáng tin --- thông tin hữu ích để quyết định đầu
tư cho khâu nào.

\textbf{Về thực tiễn}, thói quen theo từng người là lớp nền cho các ứng dụng chăm sóc người cao
tuổi trong những ngôi nhà có nhiều hơn một người --- trường hợp phổ biến nhưng ít được nghiên cứu.
Chỉ dùng cảm biến môi trường giúp giải pháp rẻ và tôn trọng quyền riêng tư.

\textbf{Giới hạn cần nói rõ}: dữ liệu thu ở Mỹ giai đoạn 2009--2011 với tối đa hai người; bất
thường trong RQ3 là bất thường \emph{cài vào}; khóa luận không khẳng định hệ thống dùng được ngay
cho nhà thật ở Việt Nam.

\section{Cấu trúc khóa luận}

\begin{itemize}
  \item \textbf{Chương~\ref{chap:background}} tổng quan nhà thông minh, nhận dạng hoạt động, bài
  toán nhà nhiều người, khai phá mẫu sinh hoạt và các kiến thức nền; định vị khoảng trống.
  \item \textbf{Chương~\ref{chap:method}} phát biểu hình thức bài toán, định nghĩa thói quen và
  trình bày toàn bộ quy trình cùng cách đánh giá.
  \item \textbf{Chương~\ref{chap:exp}} trình bày dữ liệu, thiết lập và kết quả thực nghiệm cho ba
  câu hỏi nghiên cứu, ablation và phân tích độ nhạy.
  \item \textbf{Chương~\ref{chap:conclusion}} tổng kết, nêu hạn chế và hướng phát triển.
\end{itemize}
